\section{层自适应数据流与跨 Pass 调度}
\label{sec:scheduling}

\subsection{层自适应 tile 选择}

空间 tile 高度同时影响五类容量与效率：PSUM 像素数、materialized entry、
packed-HWC 重排、line words 和 dispatch 数。对候选高度 $h_t$，本文使用以下
硬约束：
\begin{align}
W_o h_t &\le 1024,\\
W_o h_t\left\lceil K/18\right\rceil&\le32768,\\
W_o h'_t\left\lceil C_{out}/32\right\rceil&\le4096,\\
\left\lceil C_{in}/4\right\rceil\left\lceil W_i/2\right\rceil&\le2048,
\end{align}
其中 $h'_t$ 是后池化输出高度。满足容量后，选择尽可能减少 tile 数、同时不
放大参数重复包的高度。当前固定安全表为
\code{2,4,8,8,13,13,8,13,13,13,6,13,13}，分别对应 13 个卷积。

该表体现三种层策略。m0--m6 的宽特征图以较小高度限制像素和行缓存；m8、
m10、m14--m16 允许半幅或整幅高度，减少层 dispatch；m13 受长 $K$ 的
materialized cache 限制使用 8；m19 使用 6 是最紧容量点。两个检测头用
13 行，以 P4 两个 tile、P5 一个 tile 完成 255 通道输出。所有表项由同一
model spec 生成参数包、Vitis header 和论文表格，避免软件与论文各维护一份
“看似相同”的常量。

容量约束只是 tile 选择的必要条件，不是充分的性能目标。若候选高度集合为
$\mathcal{H}_l$，可以把层 $l$ 的代价写为
\begin{equation}
J_l(h)=\alpha N_{tile}(h)+\beta B_{param}(h)+\gamma B_{ifm}(h)
+\delta R_{route}(h),
\end{equation}
其中前三项分别代表 dispatch 数、按 tile 重复的参数字节和输入搬运，
$R_{route}$ 是缓存规模与跨区连线的物理风险。本文安全表没有声称求得全局
解析最优，而是在硬容量、参数窗口和最终 200 MHz 可布线条件下选择离散点。
例如 m19 从 6 增至 7 会立即违反 materialized 深度；m13 即使容量允许更大
高度，也可能因 URAM 位置和长 $K$ 控制扇出损害路由。最终选择因此必须同时
通过软件包容量检查、综合资源、place 拥塞和 route timing，不能只由 Python
层形状脚本决定。

tile 表还影响参数包。当前格式为了让每个 dispatch 自包含，会按空间 tile
重复相应 bias/weight section。减小 tile 虽降低片上容量，却会扩大参数包和
启动次数；增大 tile 则相反。安全表下 bias 与 weight 分别为 64,256 B 和
18,614,016 B，能够进入既有 64 KiB/20 MiB 窗口。生成器对总量、64 B 对齐、
DMA 长度和每层 section 完整性 fail-close，因此“更快的 tile”若导致包越界，
不会在板端运行到一半才暴露地址覆盖。

\subsection{Materialize-once/replay-many}

一个 spatial tile 的输入首先按 HWC 顺序 DMA 到物化器。物化器生成
$N_K$ 组 18-lane 窗口向量并写入 cache；随后第一个 Cout block 的 pass0
开始计算。对同一 Cout block 的 pass1--$N_K-1$，feeder 从 cache 重放对应
向量；切换到下一个 Cout block 时，输入向量仍可重放，不再读取 DDR。
因此输入 DDR 字节与 $N_C N_K$ 解耦，近似由 tile 原始 HWC 大小决定。

若原始 tile 字节数为 $B_X$，展开后的 18-lane 向量总量为 $B_M$，逐 pass
从 DDR 读取的基线输入流量近似为 $N_C N_K B_X$，物化重放则为
$B_X+B_M$ 的片外输入，其中 $B_M$ 只在片上写读。二者的收益随 $N_C N_K$
增长，但不是所有层都同样明显：m0 的 $N_K$ 很小，主要价值来自直接 HWC
窗口化；m13 的 $N_C=32,N_K=256$，若逐 pass 搬运将完全淹没计算；m19 同时
具有 5 个空间 tile 和 192 个 pass，调度与容量必须共同考虑。论文消融将记录
实际 DMA 字节而不只使用此上界模型，因为 AXIS packet、边界 padding 和
尾 tile 会使字节数偏离理想公式。

物化本身也不是免费步骤。首个 pass 之前存在填充延迟，窗口生成需要行缓冲，
并且新的 tile 只有在不覆盖活跃 reader 时才能提前进入。LASA 的目标不是让
物化时间消失，而是把它从“每个 pass 重复”变为“每个 tile 一次”，并在容量
允许时与上一上下文尾部计算重叠。对小 $N_K$ 层，若重叠收益不足，计数器会
显示 materializer 等待占比，层自适应调度可以选择更直接的原生 1$\times$1
路径，而不是强制所有层经过相同深缓存流程。

cache 采用带 epoch 的双上下文。写者完成 row-start/row-complete 后发布
静态准入摘要；读者按 active row、spec tag、minimum-needed cursor 动态检查。
当下一 tile 与当前读者的行集合无冲突时可以提前物化；否则 defer，而不是
覆盖仍会被重放的数据。该机制比固定 ping-pong 更灵活：早层大 tile 可能
占用不同 row 范围，尾 tile 也不必等待完整另一半缓存被释放。

\subsection{Weight preload 与上下文身份}

每个计算上下文由 layer/tile/Cout/pass、预期像素数、量化域、输出地址和
context epoch 唯一标识。weight preloader 在当前 bank 计算时接收下一上下文
的 18$\times$32 权重 tile，完成后置 ready/tag。计算控制器只有在
\begin{equation}
ready_{next}\land tag_{next}=tag_{issue}\land owner_{next}=free
\end{equation}
时才能 fast handoff。bias 和检测头尾 lane 的有效位与同一 context 绑定，
防止“权重是下一 pass、valid 仍是上一 Cout block”的错位。

在层与层之间，软件仍发起新的 dispatch，但 layer 内的 pass/Cout context
由硬件持续推进。这样控制寄存器写和 DMA 建链不会进入每个 $K$-pass 的内层
循环。telemetry 记录 preload miss、handoff、unexpected context、feeder
等待和所有权冲突，用于将性能归因到具体调度机制。

\subsection{Continuous PSUM 与解耦 drain}

传统反馈在每个 pass 末等待所有 PSUM 完全写入，再启动下一 pass。\LASA
允许 drain 和 feeder 对不同 packet 并行工作：当像素 $p$ 的 $P_t$ 已提交到
反馈 bank 且下一 pass 的权重/输入已就绪，$P_{t+1}$ 可在其他像素仍 drain
时开始。packet 顺序固定为先像素、后 32 通道块；read/write cursor 的距离
和 bank owner 共同保证读者不会越过写者。

对于 final pass，结果进入 requant path，不再占用反馈 bank。final drain 的
完成事件同时释放当前 PSUM owner，并允许下一上下文获取 bank。drain ring
将网络级 descriptor ready 与宽 payload 隔离后，packet/hold 数据只由注册的
FIFO return 写入；该结构既保留满队列同拍 pop+return 的 1 packet/clock，
又避免高扇出控制锥破坏路由时序。

\subsection{Fast context handoff}

如果当前上下文最后一个 token、下一权重 bank 和输入重放均已就绪，控制器
无需回到 idle 即可切换 context tag。handoff 需要额外处理同拍 alloc、release
和 done：新的 owner permit 在 context admit 或 alloc-fire 同拍通过旁路生效，
handoff-done guard 则禁止旧 done 被新上下文误认。release 是单调事件，一旦
上下文资源释放，不会因迟到状态重新占有。

本文使用三条不变量描述调度正确性：
\begin{enumerate}[leftmargin=2em]
  \item \textbf{唯一所有权：}任一 weight/PSUM/materialized bank 在一个周期
  至多属于一个未退休 context；
  \item \textbf{不覆盖：}写 cursor 不得跨越任一活跃 reader 的 minimum-needed
  row 或未消费 PSUM；
  \item \textbf{严格提交：}外部 OFM byte 顺序只由 layer、tile、pixel 和 Cout
  的字典序决定，预载和内部完成次序不得改变可见输出。
\end{enumerate}
随机反压、FIFO 空洞、满队列同拍 pop/return、异常 final pop、reset/idle/start
优先级和连续多 context 测试均围绕这些不变量构造。受控消融将分别关闭
preload、handoff 和 continuous PSUM，但每个变体首先必须通过同一 byte-exact
测试；不正确的“更快”变体不进入性能表。

三条不变量分别约束空间、时间和外部可见顺序。唯一所有权排除两个上下文在
同拍把同一 bank 当作可写资源；不覆盖把 cache 行生命周期从“固定双缓冲
半区”推广为 reader 集合；严格提交则屏蔽内部乱序，使软件无需理解预载或
drain 的完成先后。验证时不能只在理想 ready=1 下检查这些性质，因为大多数
覆盖错误出现在 return、pop、alloc、release 和 reset 同拍。测试平台因此在
每个合法边界随机插入反压，并以软件 oracle 逐周期比较 owner、cursor、count、
tag 和输出字节。系统层再以 22 节点 hash 验证局部不变量没有在长序列中遗漏。

需要强调的是，本文目前提供的是由结构不变量、定向随机测试和板端证据组成
的工程证明，并非对全部状态空间的形式化定理证明。未来可以将 bank owner、
cursor 单调性和提交次序写成 SVA，在有界 context 数下做模型检查；但在完成
之前，论文不会用“形式化证明无死锁”替代实际验证。当前可陈述的结论是：在
已覆盖的随机反压、边界和完整网络序列中未观察到所有权冲突、数据覆盖或
顺序错误，且 fail-close 门禁能把协议异常与成功完成区分开。

\subsection{计数器与利用率定义}

硬件暴露 contexts、compute-fire、IFM/OFM bytes、packet/beat、feeder wait、
context-PSUM gap、drain wait、bias/weight wait 和 unclassified cycle。对模型
总 MAC 数 $M$、PL 时间 $T_{PL}$ 和理论峰值 $P_{peak}=576f$，定义
\begin{align}
P_{eff}&=2M/T_{PL},\\
U_{array}&=P_{eff}/(2P_{peak}),\\
U_{compute}&=N_{fire}/N_{busy}.
\end{align}
前者用模型运算量和板端 wall/cycle 测量，包含 pass 尾 lane、Cout 尾块和
调度空泡；后者只观察控制器 fire，用于区分阵列内部占用与系统有效吞吐。
消融必须同时报告加速、DDR 字节、compute idle、context gap 和资源增量，
避免把减少某一计数器误写成端到端收益。
